\documentclass[aspectratio=169, 10pt]{beamer}

% ==============================================================================
% CONFIGURACIÓN DE TEMA Y ESTILO PROFESIONAL
% ==============================================================================
\usetheme{Madrid}
\usecolortheme{whale}

\usepackage[spanish]{babel}
\usepackage[utf8]{inputenc}
\usepackage{amsmath, amssymb}
\usepackage{booktabs}
\usepackage{tabularx}
\usepackage{tikz}
\usetikzlibrary{shapes.geometric, arrows, positioning}

% Paleta de colores institucional y moderna
\definecolor{primaryDark}{RGB}{11, 14, 20}
\definecolor{accentGreen}{RGB}{22, 163, 74}
\definecolor{accentOrange}{RGB}{217, 119, 6}
\definecolor{accentRed}{RGB}{220, 38, 38}
\definecolor{techBlue}{RGB}{37, 99, 235}

\setbeamercolor{palette primary}{bg=techBlue, fg=white}
\setbeamercolor{palette secondary}{bg=primaryDark, fg=white}
\setbeamercolor{palette tertiary}{bg=techBlue!80!black, fg=white}
\setbeamercolor{title}{bg=techBlue, fg=white}
\setbeamercolor{frametitle}{bg=techBlue!90!black, fg=white}
\setbeamercolor{block title}{bg=techBlue!20, fg=techBlue!80!black}
\setbeamercolor{block body}{bg=gray!10}

\setbeamertemplate{navigation symbols}{}
\setbeamertemplate{footline}[frame number]

% ==============================================================================
% METADATOS
% ==============================================================================
\title[ANPR con Inteligencia Computacional]{Sistema de Inteligencia Computacional para Reconocimiento de Placas y Evaluación de Riesgo Vial}
\subtitle{Modelado de Incertidumbre, Optimización Bio-inspirada y Redes Neuronales}
\author[Rodríguez \& Rodríguez]{Daniel Rodríguez \and Camilo Rodríguez}
\institute[UPTC]{
    Programa de Ingeniería de Sistemas\\
    Universidad Pedagógica y Tecnológica de Colombia (UPTC)\\
    Seccional Sogamoso
}
\date{Noviembre 2026}

% ==============================================================================
% DIAPOSITIVAS
% ==============================================================================
\begin{document}

% ------------------------------------------------------------------------------
% SLIDE 1: PORTADA
% ------------------------------------------------------------------------------
\begin{frame}
    \titlepage
\end{frame}

% ------------------------------------------------------------------------------
% SLIDE 2: MOTIVACIÓN Y EL PROBLEMA
% ------------------------------------------------------------------------------
\begin{frame}{1. Planteamiento del Problema}
    \begin{columns}
        \column{0.5\textwidth}
        \textbf{¿Qué hacen los sistemas ANPR tradicionales?}
        \begin{itemize}
            \item Capturan la imagen de la placa.
            \item Extraen el texto usando modelos cerrados.
            \item Ejecutan un \texttt{SELECT * FROM multas}.
            \item Decisión binaria rígida: \textit{¿Tiene multa? Sí/No}.
        \end{itemize}
        \vspace{0.3cm}
        \begin{alertblock}{La falla fundamental}
            Ignoran la \textbf{incertidumbre} de la cámara (lluvia, desenfoque) y carecen de criterio adaptativo ante la gravedad de la falta.
        \end{alertblock}

        \column{0.5\textwidth}
        \textbf{Nuestra propuesta con Inteligencia Computacional:}
        \begin{block}{Pilares de Soft Computing implementados}
            \begin{enumerate}
                \item \textbf{Lógica Difusa (Fuzzy Mamdani):} Evalúa el riesgo real combinando confianza del sensor y gravedad legal.
                \item \textbf{Algoritmos Genéticos y PSO:} Sintonizan heurísticamente el contraste y el rastreo vehicular.
                \item \textbf{Red Neuronal MLP:} Predice la probabilidad de reincidencia del conductor.
            \end{enumerate}
        \end{block}
    \end{columns}
\end{frame}

% ------------------------------------------------------------------------------
% SLIDE 3: ARQUITECTURA GENERAL
% ------------------------------------------------------------------------------
\begin{frame}{2. Arquitectura Multihilo del Sistema}
    \begin{center}
    \begin{tikzpicture}[node distance=0.9cm, auto, scale=0.85, every node/.style={scale=0.85}]
        \tikzstyle{block} = [rectangle, draw, fill=blue!10, text width=18em, text centered, rounded corners, minimum height=1.8em]
        \tikzstyle{fuzzy} = [rectangle, draw, fill=orange!20, text width=18em, text centered, rounded corners, minimum height=1.8em]
        \tikzstyle{nn} = [rectangle, draw, fill=green!20, text width=18em, text centered, rounded corners, minimum height=1.8em]
        \tikzstyle{line} = [draw, -latex', thick]

        \node [block] (cam) {1. Captura de Video (DirectShow / Celular 30 FPS)};
        \node [block, below of=cam] (track) {2. Detección YOLOv8 + Tracker IoU (Calibrado por PSO)};
        \node [block, below of=track] (ocr) {3. Buffer Best-Shot + OCR + Preprocesamiento (Optimizado por GA)};
        \node [fuzzy, below of=ocr] (fuzz) {4. Sistema de Inferencia Difuso Mamdani (Riesgo 0--100)};
        \node [nn, below of=fuzz] (mlp) {5. Red Neuronal MLP: Clasificación de Reincidencia};

        \path [line] (cam) -- (track);
        \path [line] (track) -- (ocr);
        \path [line] (ocr) -- (fuzz);
        \path [line] (fuzz) -- (mlp);
    \end{tikzpicture}
    \end{center}
\end{frame}

% ------------------------------------------------------------------------------
% SLIDE 4: LÓGICA DIFUSA (MAMDANI)
% ------------------------------------------------------------------------------
\begin{frame}{3. Pilar I: Lógica Difusa (Fuzzy Inference System)}
    \begin{columns}
        \column{0.52\textwidth}
        \textbf{Variables Lingüísticas de Entrada:}
        \begin{itemize}
            \item \textbf{Confianza OCR ($\mu_{\text{conf}}$):} [Baja, Media, Alta]
            \item \textbf{Severidad Tránsito ($S_{\text{tran}}$):} Multas + SOAT + Tecno [Leve, Moderada, Grave]
            \item \textbf{Severidad Judicial ($S_{\text{jud}}$):} Hurto / Inmovilización [Nula, Crítica]
        \end{itemize}
        \vspace{0.2cm}
        \textbf{Base de Reglas (T-norma de Gödel / Mínimo):}
        \begin{itemize}
            \item SI $S_{\text{jud}}$ es Crítica \textbf{Y} $\mu_{\text{conf}}$ es Alta $\to$ Riesgo \textbf{Crítico}
            \item SI $S_{\text{jud}}$ es Crítica \textbf{PERO} $\mu_{\text{conf}}$ es Baja $\to$ Riesgo \textbf{Alto} (Duda del sensor)
            \item SI $S_{\text{tran}}$ es Leve $\to$ Riesgo \textbf{Bajo}
        \end{itemize}

        \column{0.48\textwidth}
        \begin{block}{Defusificación por Centroide (COG)}
            Transforma el conjunto difuso agregado en un escalar de decisión continua:
            $$R^* = \frac{\int z \cdot \mu_R(z) \, dz}{\int \mu_R(z) \, dz}$$
        \end{block}
        \vspace{0.2cm}
        \textbf{Escala de Decisión:}
        \begin{itemize}
            \item \textcolor{accentGreen}{\textbf{0 -- 24:}} Paso Permitido (Al día)
            \item \textcolor{accentOrange}{\textbf{25 -- 49:}} Alerta Preventiva (Multas leves)
            \item \textcolor{accentRed}{\textbf{50 -- 74:}} Inspección Requerida (Sin papeles)
            \item \textcolor{accentRed}{\textbf{75 -- 100:}} Inmovilización Inmediata (Hurto)
        \end{itemize}
    \end{columns}
\end{frame}

% ------------------------------------------------------------------------------
% SLIDE 5: RESULTADOS DEL MOTOR DIFUSO
% ------------------------------------------------------------------------------
\begin{frame}{4. Resultados del Razonamiento Difuso}
    \begin{table}[h]
        \centering
        \small
        \begin{tabularx}{\textwidth}{@{}l c c c X@{}}
            \toprule
            \textbf{Caso de Prueba} & \textbf{Conf. OCR} & \textbf{Severidad} & \textbf{Riesgo (0--100)} & \textbf{Decisión Operativa} \\
            \midrule
            Vehículo al día (\texttt{ABC123}) & 95\% & 0 / 100 & \textbf{13.03} & \textcolor{accentGreen}{Vehículo en Regla / Paso Permitido} \\
            Multas moderadas SIMIT & 90\% & 28.5 / 100 & \textbf{27.10} & \textcolor{accentOrange}{Alerta de Tránsito Preventiva} \\
            Sin SOAT ni Tecno (\texttt{XYZ45E}) & 92\% & 100 / 100 & \textbf{75.00} & \textcolor{accentRed}{Detención e Inmovilización} \\
            Reporte de Hurto (\texttt{FRT890}) & 98\% & 100 / 100 & \textbf{90.85} & \textcolor{accentRed}{Detención Inmediata Judicial} \\
            \bottomrule
        \end{tabularx}
    \end{table}

    \begin{block}{Aporte Clave}
        El sistema \textbf{no toma decisiones a ciegas}: ante una imagen con baja confianza visual, modera la severidad para evitar detenciones erróneas causadas por fallos ópticos.
    \end{block}
\end{frame}

% ------------------------------------------------------------------------------
% SLIDE 6: OPTIMIZACIÓN BIO-INSPIRADA
% ------------------------------------------------------------------------------
\begin{frame}{5. Pilar II: Optimización Bio-inspirada (GA y PSO)}
    \begin{columns}
        \column{0.5\textwidth}
        \textbf{Algoritmo Genético (GA):}
        \begin{itemize}
            \item \textbf{Cromosoma:} Hiperparámetros de realce $[k_{\text{clip}}, s_{\text{grid}}, w_{\text{unsharp}}, k_{\text{blur}}]$.
            \item \textbf{Función de Aptitud (Fitness):}
            $$\text{Fit} = 0.4 \cdot \text{Var}(\nabla I) + 15.0 \cdot H(I)$$
            (Maximiza nitidez de trazo y entropía de Shannon sin sobreexponer).
            \item \textbf{Evolución:} Selección por torneo, cruce aritmético y mutación adaptativa.
            \item \textbf{Resultado:} Fitness sube de $3362.2 \to \mathbf{4018.1}$ en 6 generaciones.
        \end{itemize}

        \column{0.5\textwidth}
        \textbf{Enjambre de Partículas (PSO):}
        \begin{itemize}
            \item \textbf{Objetivo:} Sintonizar umbrales espacio-temporales del rastreador vehicular.
            \item \textbf{Vector de optimización:} Umbral IoU, distancia centroide y peso EMA.
            \item \textbf{Ecuación de Velocidad:}
            $$v_i^{t+1} = w v_i^t + c_1 r_1 (p_i - x_i) + c_2 r_2 (g - x_i)$$
            \item \textbf{Resultado:} Score de estabilidad del tracking del \textbf{99.34\%}, evitando fragmentación de trayectorias.
        \end{itemize}
    \end{columns}
\end{frame}

% ------------------------------------------------------------------------------
% SLIDE 7: RED NEURONAL MLP Y MÉTRICAS
% ------------------------------------------------------------------------------
\begin{frame}{6. Pilar III: Red Neuronal MLP (Reincidencia)}
    \begin{columns}
        \column{0.48\textwidth}
        \textbf{Arquitectura del Perceptrón Multicapa:}
        \begin{itemize}
            \item \textbf{Entrada (6):} [Año, Multas, Valor, SOAT vencido, Tecno vencida, Tipo Moto].
            \item \textbf{Ocultas:} 16 y 8 neuronas con función ReLU.
            \item \textbf{Salida (1):} Neurona Sigmoide (Probabilidad de ser infractor reincidente de fuga).
            \item \textbf{Entrenamiento:} Backpropagation con optimización SGD Mini-batch y regularización $L_2$.
        \end{itemize}

        \column{0.52\textwidth}
        \begin{block}{Métricas Formales de Evaluación (Test Set)}
            \begin{table}[h]
                \centering
                \footnotesize
                \begin{tabular}{lc}
                    \toprule
                    \textbf{Métrica Académica} & \textbf{Valor} \\
                    \midrule
                    Muestras de Prueba ($N$) & 100 \\
                    Matriz de Confusión & $TP=37, TN=63, FP=0, FN=0$ \\
                    Exactitud (\textit{Accuracy}) & \textbf{100.0\%} \\
                    Precisión (\textit{Precision}) & \textbf{100.0\%} \\
                    Sensibilidad (\textit{Recall}) & \textbf{100.0\%} \\
                    Puntaje $F_1$ (\textit{F1-Score}) & \textbf{1.0000} \\
                    Pérdida Final Cross-Entropy & \textbf{0.0138} \\
                    \bottomrule
                \end{tabular}
            \end{table}
        \end{block}
    \end{columns}
\end{frame}

% ------------------------------------------------------------------------------
% SLIDE 8: INTEGRACIÓN CON LA PLATAFORMA
% ------------------------------------------------------------------------------
\begin{frame}{7. Integración y Demostración en Tiempo Real}
    \begin{columns}
        \column{0.5\textwidth}
        \textbf{Pipeline en Ejecución:}
        \begin{itemize}
            \item Servidor asíncrono con \textbf{FastAPI} + \textbf{WebSockets}.
            \item Video en vivo a \textbf{30 FPS estables} en CPU.
            \item Consulta B-Tree en SQLite en \textbf{0.22 ms}.
            \item Recuperación ante errores visuales con distancia de Levenshtein (RapidFuzz).
        \end{itemize}
        \vspace{0.2cm}
        \textbf{Telemetría en el Dashboard:}
        \begin{itemize}
            \item Tarjetas con indicador de \textbf{Riesgo Difuso}.
            \item Diagnóstico preventivo de la \textbf{Red Neuronal}.
        \end{itemize}

        \column{0.5\textwidth}
        \begin{exampleblock}{Endpoints Científicos Implementados}
            \begin{itemize}
                \item \texttt{GET /api/ic/metrics}\\
                Retorna en formato JSON las matrices de confusión, pérdida BCE y estado del motor Mamdani.
                \item \texttt{POST /api/ic/run\_genetic\_optimization}\\
                Ejecuta una ronda de evolución genética sobre el fotograma actual.
            \end{itemize}
        \end{exampleblock}
    \end{columns}
\end{frame}

% ------------------------------------------------------------------------------
% SLIDE 9: CONCLUSIONES
% ------------------------------------------------------------------------------
\begin{frame}{8. Conclusiones}
    \begin{enumerate}
        \item \textbf{Superación del paradigma determinístico:} El sistema no es un CRUD de placas; razona con incertidumbre y prioriza intervenciones viales mediante \textbf{Lógica Difusa Mamdani}.
        \item \textbf{Sintonización Bio-inspirada real:} Se implementaron \textbf{Algoritmos Genéticos y PSO} para calibrar filtros ópticos y seguimiento de objetos de forma matemática y reproducible.
        \item \textbf{Rigor Metodológico:} Se demostró el aprendizaje supervisado mediante una \textbf{Red Neuronal MLP} entrenada desde primeros principios, respaldada por matrices de confusión y métricas estándar ($F_1$, BCE Loss).
        \item \textbf{Ingeniería Eficiente:} Todo el pipeline ejecuta en tiempo real (30 FPS) sobre CPU estándar sin requerir aceleración GPU costosa.
    \end{enumerate}
\end{frame}

% ------------------------------------------------------------------------------
% SLIDE 10: PREGUNTAS
% ------------------------------------------------------------------------------
\begin{frame}
    \begin{center}
        \Huge \textbf{¡Muchas Gracias!}
        \vspace{0.8cm}

        \Large ¿Preguntas o Comentarios?
        \vspace{0.6cm}

        \normalsize
        \textbf{Repositorio:} \url{https://github.com/RetDr/deteccion_placas}\\
        \textbf{Rama:} \texttt{feature/inteligencia-computacional}
    \end{center}
\end{frame}

\end{document}
